\chapter{数值稳定性与 Python 互操作}
\index{numerical stability}\index{floating point}\index{pybind11}\index{GIL}\index{NumPy}

\chapteropening{能解释浮点比较边界，使用稳定求和与在线方差，区分金额和收益率的表示责任，并设计明确 dtype、布局、GIL 与所有权的批量绑定。}{一个无第三方依赖的数值核心、可选 pybind11 模块及 Python/C++ 一致性证据。}{第 5、12、14 章；会维护类不变量、写行为测试并以独立基准约束测量。}

\section{问题：先验证算法，再跨边界}

Python 研究代码移到 C++ 时，两个问题常被错误地混在一起：数值结果为何改变，以及调用为什么没有变快。\path{code/ch16} 把它们拆成一个无第三方依赖的 C++20 核心和一个可选 pybind11 薄绑定。核心始终可构建；只有环境已经提供 pybind11 时才生成 Python 模块。

\begin{pythonbridge}
NumPy 数组表达式和 C++ 批量 kernel 都能把逐元素循环移出 Python；真正的边界由 dtype、布局、所有者和错误政策共同决定。
\pythoncompare{NumPy 用数组与 ufunc；C++ 暴露接收连续批次的静态类型函数。}{Python 数组通过所有者或 base 对象保活；C++ span 只借用，返回数组必须拥有结果。}{NumPy 在调用时检查 dtype/shape；C++ 核心接口先编译，绑定再验证运行期元数据。}{Python 侧通常抛 TypeError/ValueError；未验证的 C++ 借用可能悬空或越界。}
\end{pythonbridge}

从项目根目录运行独立快照：

\begin{lstlisting}[language=bash,numbers=none]
cmake -S code/ch16 -B build/ch16-debug -DCMAKE_BUILD_TYPE=Debug
cmake --build build/ch16-debug -j
ctest --test-dir build/ch16-debug --output-on-failure
./build/ch16-debug/ch16_stability_report
./build/ch16-debug/ch16_boundary_report
\end{lstlisting}

首个稳定输出为：

\begin{lstlisting}[numbers=none]
stability-ok exact-0.1+0.2=false tolerant=true
  naive=3.0 stable=4.0 oracle=4.0
tolerance-source input-abs=0.000e+00 operations=2
  scale=3.000e-01 abs=1.332e-16 rel=4.441e-16
\end{lstlisting}

独立手算基准是 $10^{16}+1-10^{16}+3=4$。程序若只打印“计算完成”或拿同一个函数生成 expected，不能证明算法正确。

\section{浮点表示、比较与容差来源}

\texttt{double} 通常按 IEEE 754 binary64 保存符号、指数和有限位有效数。$0.1$ 与 $0.2$ 在二进制中都是无限小数，字面量先舍入到相邻可表示值，再参与加法；因此 \texttt{0.1 + 0.2 == 0.3} 为假并不表示 CPU 随机算错。\texttt{==} 仍适合整数、枚举，以及本来就要求位级完全相同的协议值；问题是把近似实数误当成精确编码。

\path{numerics.cpp} 的比较规则同时给绝对与相对下界：

\begin{lstlisting}[language=C++,numbers=none]
const double difference = std::abs(lhs - rhs);
const double scale = std::max(std::abs(lhs), std::abs(rhs));
return difference <=
       std::max(absolute_tolerance, relative_tolerance * scale);
\end{lstlisting}

绝对容差保护接近零的值；相对容差随量级缩放。可运行的 \texttt{tolerance\_from\_budget} 接收输入绝对误差、预计舍入操作数和结果尺度。本例输入误差为 0、加法含两次舍入、尺度为 0.3，按 \texttt{double} epsilon 导出上面的绝对与相对界限。报价精度或参考实现已有误差时，应计入 input-abs；容差过大仍会吞掉真实回归。

Python 的 \texttt{math.isclose} 也使用相对与绝对容差，但默认值、NaN 规则和调用点不会自动成为 C++ 合约。失败可由直接断言 \texttt{simple\_return(100, 102.5) == 0.025} 触发；诊断是值只在末位不同，根因是两次舍入，修复是给出业务来源的容差并保留独立 expected。

\section{求和顺序、补偿项与在线方差}

浮点加法不满足实数意义下的结合律。顺序求和先执行 $10^{16}+1$ 时，较小的 1 会被舍掉；之后再减 $10^{16}$ 已无法恢复。\texttt{std::accumulate} 能表达归约，却仍按给定顺序反复做同一种加法，不能自动补偿已丢低位。误差如何累积、算法如何改善稳定性，需要结合问题尺度分析，不能只看一次输出\autocite{higham2002accuracy}。

本章用 Neumaier 补偿求和：\texttt{total} 保存主和，\texttt{correction} 收集每次加法中未进入主和的低位。当新数绝对值更大时，误差表达式的方向也随之交换。完整可运行入口如下；核心实现来自同一静态库：

\lstinputlisting[caption={朴素求和、补偿求和与独立判定基准}]{code/ch16/experiments/stability_report.cpp}

方差也不能机械套用 $E[X^2]-E[X]^2$：两项很大而差很小时会发生灾难性消减。\texttt{sample\_moments} 使用 Welford 更新，每读一个值就更新 count、mean 与平方离差和，最后明确除以 $n-1$，所以输出叫 sample variance；若业务要总体方差，则分母与名称必须同时改变。

Python 的 \texttt{math.fsum} 提供高精度求和工具，NumPy 的归约还可能按 dtype、轴和实现选择不同顺序；C++ 补偿算法不会自动等同于某个 NumPy 版本。失败触发是把输入顺序换成大数、小数、反向大数并观察朴素和偏离独立判定基准；诊断类别是 cancellation/order sensitivity，根因是有效位丢失，修复是改变算法、提高精度或重构问题，而非扩大测试容差。

\section{金额、收益率与方差需要不同表示}

金额通常有法定或业务最小单位。本章示例把已完成舍入的金额保存为 \texttt{int64\_t} cents。转换函数只在明确入口把 major units 乘 100，并用 \texttt{std::round} 定义半数远离零。输入 10.125 可由二进制精确表示，得到 1013 分；真实 CSV 若含十进制文本，应优先按字符串规则解析或使用经批准的 decimal 库，不能先丢精度再声称精确。

收益率是比值，通常保留 \texttt{double}，并记录 simple return 还是 log return、价格复权和缺失值规则。方差的单位是收益率平方，还必须声明 sample 或 population。运行 \texttt{ch16\_finance\_report} 得到：

\begin{lstlisting}[numbers=none]
numeric-policy money=int64-cents rounded-cents=1013
  return=0.025000 mean=0.006667 sample-variance=0.000633
  overflow-rejected=true
\end{lstlisting}

Python 的整数也能无损保存 cents，\texttt{decimal.Decimal} 能表达十进制规则且整数不会溢出；C++ 固定宽度整数必须检查范围，标准库也没有金融 decimal 的直接替代品。失败可由向金额入口传 NaN、无穷或超出 \texttt{int64\_t} 的值触发；稳定类别分别是 invalid input 与 overflow，根因是表示域不匹配，修复是边界拒绝而不是静默截断。

\section{逐元素调用、批量复制与零拷贝借用}

跨 Python/C++ 边界有固定成本：参数转换、动态分派、错误翻译和可能的 GIL 操作。把四个元素逐个调用 C++ 会付四次边界成本；一次批量复制只跨一次边界，但分配并复制；零拷贝借用同样只跨一次，却把布局与寿命前置条件交给接口。

\path{boundary_report.cpp} 让三条路径处理 \texttt{\{0.25,-0.125,0.5,0.375\}}，共同 checksum 为 1：

\begin{lstlisting}[numbers=none]
boundary-ok rows=4 scalar-calls=4 bulk-copies=1
  zero-copy-borrows=1 checksum=1.000000
\end{lstlisting}

调用次数和复制次数只描述成本模型，尚不足以推出耗时结论。小数组可能由转换成本主导，大数组复制后反而可能因连续布局更快；应像第 14 章那样固定结果、规模、预热和样本再测量。

Python/NumPy 的向量化把循环移到扩展内部，但“写成数组表达式”不保证零拷贝，也不保证 dtype 未转换。失败触发是对每个行情点调用一次绑定后只测 C++ 循环；诊断是 profiler 显示大量 binding/call overhead，根因是边界粒度太细，修复是设计批量领域操作。

\section{dtype、连续性与缓冲所有权}

NumPy 数组由地址和元数据共同定义。dtype 决定一个元素如何解释；shape 给每维长度；stride 给同一维相邻元素的字节距离。切片 \texttt{values[::2]} 可以共享原数组却不连续。若 C++ 把 float32 地址强转成 \texttt{double*}，每八字节会拼接两个四字节元素；若忽略 stride，读取的也不是逻辑相邻值。

本章的 \texttt{BatchView} 要求一维 float64 且 stride 等于 \texttt{sizeof(double)}，不允许隐式 forcecast。错误输入得到稳定边界：

\begin{lstlisting}[numbers=none]
dtype-error expected=float64 actual=float32
layout-error expected=c-contiguous stride-bytes=16
ownership-error temporary-view-rejected owner-missing=true
\end{lstlisting}

严格拒绝适合教学和延迟敏感路径；另一种合法 API 可以明确命名 \texttt{allow\_copy}，由它转换 dtype 与布局并报告发生了复制。两种语义不能藏在同一个“看起来零拷贝”的函数里。

Python 视图常通过 \texttt{base} 保持源对象存活；C++ \texttt{span} 只是指针加长度，不拥有元素。故意失败的 \path{temporary_buffer_view.cpp} 返回局部 vector 的 span，离开函数即悬空；ASan 稳定报告 \texttt{heap-use-after-free}。pybind11 返回借用数组时必须把原 Python 对象设为 base，返回新 C++ 缓冲时则用 capsule 或拥有型数组托管释放。

\section{GIL 与可选 pybind11 薄绑定}

Python 调入扩展时当前线程持有 GIL。长时间纯 C++ 循环若一直持有它，会阻止其他 Python 线程执行 Python 字节码；但释放 GIL 只移除解释器锁，不会替 C++ 共享状态建立 mutex 或 happens-before。pybind11 提供作用域式释放工具，同时要求释放期间不调用需要 GIL 的 Python API\autocite{pybind11docs}。

\path{bindings/module.cpp} 先在持有 GIL 时取得 \texttt{buffer\_info}，随后只在不调用 Python API 的核心求和或填充循环周围创建 \texttt{py::gil\_scoped\_release}；作用域结束自动重新取得 GIL，之后才构造 Python 异常或返回对象。若循环中需要回调 Python、调整引用计数或创建数组，就不能在释放区间执行这些操作。

可选构建遵守依赖边界：\texttt{find\_package(pybind11 CONFIG QUIET)} 找不到依赖时只打印状态并继续构建核心；找到时生成 \texttt{quant\_ch16}，其 \texttt{sum\_returns} 仍调用同一 \texttt{sum\_batch}。Python 验收脚本检查 float64 连续输入、错误 dtype、切片 stride，以及删除原名称后借用视图仍由 base 保活。

Python 的线程同样受 GIL 约束，但 NumPy 或其他扩展是否释放、释放多久属于各函数实现契约。失败触发是两个 Python 线程调用长 C++ 循环却完全串行；profiler 显示第二线程没有 Python 进展，根因通常是持锁计算。修复是在纯 C++ 且对象寿命已固定的最小区间释放，并用 ThreadSanitizer/行为测试另证 C++ 线程安全。

无需 pybind11 的 \path{gil_probe.cpp} 还用 CPython 扩展 API 建立可运行诊断：新 C++ 线程必须取得 GIL 后才设置原子进度。调用方持有 GIL 等待时进度稳定为 false；仅在等待区间释放后才为 true，测试输出 \texttt{gil-diagnostic held-progress=false released-progress=true}。探针以 deadline 防止故障无限挂起，但不靠 sleep 猜测调度。

\section{章末三层练习}

\subsection*{随堂验证汇总}

\begin{enumerate}
  \item 只把 input-abs 从 0 改为 $10^{-4}$；输出 absolute 应约为 $10^{-4}$，并解释这代表输入精度变化而非“为了过测试”。
  \item 只把数组顺序改为 \texttt{\{1e16, -1e16, 1, 3\}}；朴素和应变为 4，稳定和仍为 4，并解释“这次朴素正确”为何不能证明算法稳健。
  \item 只把金额改为 $-10.125$，按本章口径应得到 $-1013$ 分；再写一句说明退款是否允许为负属于领域规则而非浮点规则。
  \item 只把输入行数从 4 改为 8；三条 checksum 仍须相同，scalar-calls 应为 8，而 bulk-copies 与 zero-copy-borrows 仍各为 1。
  \item 只把步长 16 改回 8，layout 错误应消失；dtype 仍为 float32 时必须继续拒绝，说明两个前置条件彼此独立。
  \item 只把 \texttt{gil\_scoped\_release} 移出数据指针取得之前，指出为什么此时访问 \texttt{py::array} 元数据不再安全；正确答案不是“多加一个锁”，而是缩回释放作用域。
\end{enumerate}

复现层固定数值与边界输出；修改层改变批大小或缺失值分布并保持独立判定基准；开放层为一段 Python 热点设计 dtype、所有权、GIL 与测量契约。

输出任务：提交 \path{code/ch16} 的 clean-build 命令、容差预算与三组精确输出，以及一页边界说明。说明须给出金额舍入口径、收益与方差定义；比较逐元素、批量复制和零拷贝；列出 dtype、stride、所有者与 GIL 的前置条件。附 ASan 的 \texttt{heap-use-after-free} 与 GIL probe 两类稳定证据。环境有 pybind11 时再提交可选模块输出；没有时提交“dependency-free core only”配置证据。

自测：
\begin{enumerate}
  \item 为什么 \texttt{0.1 + 0.2 == 0.3} 失败，而整数 cents 可以精确比较？
  \item 绝对容差与相对容差分别保护什么量级，容差应从哪里来？
  \item Neumaier 补偿项保存了什么，为什么不能只把容差调大？
  \item Welford 方差的 count、mean 与平方离差和如何逐步变化？
  \item 金额、收益率和样本方差为何不应共用一种表示口径？
  \item 批量复制与零拷贝各把什么成本和责任交给调用者？
  \item dtype、shape、stride 与所有者各阻止哪类错误？
  \item 释放 GIL 能解决什么，又绝不能替代什么？
\end{enumerate}


\section{本章小结}

数值可靠性来自表示、算法和验收口径；互操作性能来自足够粗的边界。先以独立判定基准验证无依赖核心，再把 dtype、连续性、所有者和 GIL 写进薄绑定协议。
